\documentclass[conference]{IEEEtran}
\usepackage{booktabs}
\usepackage{hyperref}
\usepackage{graphicx}
\usepackage{xcolor}
\usepackage{url}

\title{Does Project-Specific In-Context Learning Improve the Reliability of LLM-Generated Code Summaries?}

% TODO: Replace with actual author name(s) and affiliation(s)
\author{
  \IEEEauthorblockN{Author Name}
  \IEEEauthorblockA{Affiliation\\Email}
}

\begin{document}

\maketitle

% ============================================================
\begin{abstract}
Large language models (LLMs) can generate natural-language summaries of source
code, but generic prompting often produces summaries that are vague or
inaccurate. We investigate whether \emph{project-specific in-context learning}
(ICL)---providing the LLM with example function--summary pairs drawn from the
same project---improves summary reliability. In a small-scale exploratory study
covering 15~Python functions from three open-source repositories (Requests,
pandas, scikit-learn), we compare generic prompting against project-specific
ICL and manually label each generated summary as \textsc{Correct},
\textsc{Partial}, or \textsc{Incorrect}. Project-specific ICL achieved a 73.3\%
correctness rate compared to 46.7\% for generic prompting, and reduced the
incorrect-summary rate from 26.7\% to 6.7\%. While the sample size limits
statistical generalizability, the results suggest that grounding LLM prompts in
project-specific context is a promising direction for more reliable automated
code summarization. All data, prompts, outputs, evaluation labels, and scripts
are publicly available for replication.
\end{abstract}

% ============================================================
\section{Introduction}
\label{sec:introduction}

Automated code summarization aims to generate concise natural-language
descriptions of source code. Recent advances in large language models (LLMs)
have made it possible to generate plausible code summaries with minimal
task-specific training~\cite{brown2020language}. However, generic prompting
approaches---where the LLM receives only a summarization instruction and the
target function---often produce summaries that are semantically inaccurate or
too generic to be useful.

Project-specific in-context learning (ICL) provides the LLM with example
function--summary pairs from the same project as the target function. The
hypothesis is that project-specific context helps the LLM align its output
with the project's conventions, vocabulary, and behavioral patterns.

\textbf{Research Question:} Does project-specific in-context learning improve
the reliability of LLM-generated code summaries compared to generic prompting?

% ============================================================
\section{Related Work}
\label{sec:related-work}

Ahmed and Devanbu~\cite{ahmed2022fewshot} explored few-shot training of LLMs
for project-specific code summarization, finding that project-specific examples
can improve summary quality. Ahmad et al.~\cite{ahmad2020neuralcodesum}
proposed NeuralCodeSum, a Transformer-based approach trained on large corpora
of code--summary pairs. Yun et al.~\cite{yun2024pcodesum} introduced
P-CodeSum, which uses in-context learning for project-specific summarization.

Our work differs in that we conduct a controlled comparison between generic
prompting and project-specific ICL on a curated set of functions, with manual
evaluation labels, and provide a complete replication package.

% ============================================================
\section{Literature Search Strategy}
\label{sec:literature-search}

\subsection{Search Protocol}
We searched Google Scholar using Publish or Perish with a broad query
(\textit{``code summarization'' OR ``source code summarization'' OR ``code
comment generation'' OR ``code documentation generation''}) and six refined
queries targeting project-specific, few-shot, ICL, and evaluation aspects.
We focused on top SE venues (ICSE, FSE, ASE, MSR, TSE, TOSEM) and related
NLP/ML venues (ACL, EMNLP, NeurIPS), primarily since 2015.

\subsection{Citation Ranking and Knee Analysis}
We collected the top 100 papers by citation count and applied a knee analysis
(maximum perpendicular distance method) to the citation curve. The knee
occurred at rank~22, corresponding to a citation threshold of~167. The 22
above-the-knee papers form our primary literature corpus.
The complete 100-paper citation list and knee analysis data are included in
the artifact repository.

\subsection{Thematic Classification}
We classified the 22 above-the-knee papers into four themes:
\textbf{P}~(Project-Specific Context, 6~papers),
\textbf{I}~(In-Context Learning and Prompting, 4~papers),
\textbf{M}~(Code Summarization Methods, 19~papers), and
\textbf{E}~(Evaluation and Reliability, 6~papers).
The full intersection $P \cap I \cap M \cap E = 0$: no existing paper
addresses all four themes simultaneously, defining the research gap our
study targets.

% ============================================================
\section{Experimental Rig}
\label{sec:experimental-rig}

\subsection{Dataset}
We selected 15~Python functions from three widely-used open-source
repositories: Requests~\cite{requests}, pandas~\cite{pandas}, and
scikit-learn~\cite{scikitlearn}. Each repository contributes 5~functions
spanning utility, data-processing, validation, and conversion categories.
Table~\ref{tab:dataset-composition} summarizes the dataset composition.

% Include generated table from outputs/latex_tables.tex
% % Table: Dataset Composition
\begin{table}[ht]
\centering
\caption{Dataset Composition by Repository}
\label{tab:dataset-composition}
\begin{tabular}{lcc}
\toprule
Repository & Functions & Proportion \\
\midrule
requests & 5 & 33.3\% \\
pandas & 5 & 33.3\% \\
scikit-learn & 5 & 33.3\% \\
\midrule
\textbf{Total} & \textbf{15} & \textbf{100.0\%} \\
\bottomrule
\end{tabular}
\end{table}



% Table: Target Functions
\begin{table}[ht]
\centering
\caption{Target Functions Used in the Study}
\label{tab:target-functions}
\begin{tabular}{llll}
\toprule
ID & Repository & Module & Function \\
\midrule
F01 & requests & \texttt{requests.utils} & \texttt{requote\_uri} \\
F02 & requests & \texttt{requests.utils} & \texttt{get\_encoding\_from\_headers} \\
F03 & requests & \texttt{requests.utils} & \texttt{select\_proxy} \\
F04 & requests & \texttt{requests.sessions} & \texttt{merge\_setting} \\
F05 & requests & \texttt{requests.models} & \texttt{request\_url} \\
F06 & pandas & \texttt{pandas.core.common} & \texttt{flatten} \\
F07 & pandas & \texttt{pandas.api.types} & \texttt{is\_list\_like} \\
F08 & pandas & \texttt{pandas.core.algorithms} & \texttt{unique} \\
F09 & pandas & \texttt{pandas.core.reshape.concat} & \texttt{concat} \\
F10 & pandas & \texttt{pandas.core.tools.datetimes} & \texttt{to\_datetime} \\
F11 & scikit-learn & \texttt{sklearn.utils.validation} & \texttt{check\_array} \\
F12 & scikit-learn & \texttt{sklearn.utils.validation} & \texttt{check\_X\_y} \\
F13 & scikit-learn & \texttt{sklearn.model\_selection} & \texttt{train\_test\_split} \\
F14 & scikit-learn & \texttt{sklearn.metrics.pairwise} & \texttt{cosine\_similarity} \\
F15 & scikit-learn & \texttt{sklearn.preprocessing} & \texttt{normalize} \\
\bottomrule
\end{tabular}
\end{table}



% Table: Result Summary
\begin{table}[ht]
\centering
\caption{Summary of Evaluation Results}
\label{tab:result-summary}
\begin{tabular}{lccccc}
\toprule
Method & Correct & Partial & Incorrect & Total & Correctness Rate \\
\midrule
Generic Prompting & 7 & 4 & 4 & 15 & 46.7\% \\
Project-Specific ICL & 11 & 3 & 1 & 15 & 73.3\% \\
\bottomrule
\end{tabular}
\end{table}



% Table: Label Percentages
\begin{table}[ht]
\centering
\caption{Label Distribution (\%) by Method}
\label{tab:label-percentages}
\begin{tabular}{lccc}
\toprule
Method & Correct (\%) & Partial (\%) & Incorrect (\%) \\
\midrule
Generic Prompting & 46.7 & 26.7 & 26.7 \\
Project-Specific ICL & 73.3 & 20.0 & 6.7 \\
\bottomrule
\end{tabular}
\end{table}



% Table: Reproducibility Baseline Status
\begin{table}[ht]
\centering
\caption{Reproducibility Status of Referenced Baseline Systems}
\label{tab:baseline-status}
\begin{tabular}{llll}
\toprule
System & Repo Available & Status & Issue \\
\midrule
Ahmed and Devanbu 2022 Few-shot project-specific code summarization & Yes & failed unmodified & Repo cloned, environment worked, script reached API call, then failed because code-davinci-002 is deprecated \\
Ahmad et al. 2020 NeuralCodeSum & Yes & failed unmodified & Repo cloned, dependencies repaired, pipeline launched, then failed because required Java dataset files were missing and download path did not produce a valid archive \\
Yun et al. 2024 P-CodeSum & Yes & partial & Repo installed successfully, scripts launched, but default train.sh/evaluate.sh point to nonexistent data paths; repo contains alternate data dirs; log files created successfully \\
\bottomrule
\end{tabular}
\end{table}



% Table: Above-the-Knee Papers
\begin{table*}[ht]
\centering
\caption{Top 22 Above-the-Knee Papers by Citation Count}
\label{tab:above-knee}
\begin{tabular}{clc}
\toprule
Rank & Title & Citations \\
\midrule
1 & Summarizing source code using a neural attention model & 1009 \\
2 & Deep code comment generation & 922 \\
3 & A transformer-based approach for source code summarization & 567 \\
4 & Improving automatic source code summarization via deep reinforcement learning & 535 \\
5 & On the use of automated text summarization techniques for summarizing source code & 527 \\
6 & Improved code summarization via a graph neural network & 428 \\
7 & Few-shot training LLMs for project-specific code-summarization & 424 \\
8 & Summarizing source code with transferred API knowledge & 392 \\
9 & Supporting program comprehension with source code summarization & 385 \\
10 & Deep code comment generation with hybrid lexical and syntactical information & 354 \\
11 & Retrieval-based neural source code summarization & 329 \\
12 & Automatic documentation generation via source code summarization of method context & 322 \\
13 & Retrieval augmented code generation and summarization & 302 \\
14 & Code generation as a dual task of code summarization & 288 \\
15 & Semantic similarity metrics for evaluating source code summarization & 271 \\
16 & Improving automated source code summarization via an eye-tracking study of programmers & 271 \\
17 & Large language models are few-shot summarizers: Multi-intent comment generation via in-context learning & 244 \\
18 & Retrieval-augmented generation for code summarization via hybrid GNN & 236 \\
19 & Automatic source code summarization of context for Java methods & 234 \\
20 & Automatic semantic augmentation of language model prompts (for code summarization) & 229 \\
21 & Evaluating source code summarization techniques: Replication and expansion & 179 \\
22 & Recommendations for datasets for source code summarization & 167 \\
\bottomrule
\end{tabular}
\end{table*}



% Table: Thematic Group Sizes
\begin{table}[ht]
\centering
\caption{Thematic Group Sizes (22 Above-the-Knee Papers)}
\label{tab:thematic-groups}
\begin{tabular}{lc}
\toprule
Theme & Papers \\
\midrule
P: Project-Specific Context & 6 \\
I: In-Context Learning \& Prompting & 4 \\
M: Code Summarization Methods & 19 \\
E: Evaluation \& Reliability & 6 \\
\bottomrule
\end{tabular}
\end{table}



% Table: Venn Disjoint Region Counts
\begin{table}[ht]
\centering
\caption{Disjoint Venn Region Counts for Thematic Classification}
\label{tab:venn-counts}
\begin{tabular}{lc}
\toprule
Region & Count \\
\midrule
Only P & 1 \\
Only I & 0 \\
Only M & 10 \\
Only E & 0 \\
P $\cap$ I only & 1 \\
P $\cap$ M only & 2 \\
P $\cap$ E only & 0 \\
I $\cap$ M only & 1 \\
I $\cap$ E only & 1 \\
M $\cap$ E only & 4 \\
P $\cap$ I $\cap$ M & 1 \\
P $\cap$ I $\cap$ E & 0 \\
P $\cap$ M $\cap$ E & 1 \\
I $\cap$ M $\cap$ E & 0 \\
P $\cap$ I $\cap$ M $\cap$ E & 0 \\
\bottomrule
\end{tabular}
\end{table}



% Table: Reproduction Study Outcomes
\begin{table*}[ht]
\centering
\caption{Baseline Reproduction Study Outcomes}
\label{tab:reproduction-study}
\begin{tabular}{lccll}
\toprule
System & Repo Found & Runs & Feasible & Decision \\
\midrule
Ahmed \& Devanbu 2022 & Yes & No & Maybe & Conceptual baseline; API deprecated \\
Ahmad et al.\ 2020 (NeuralCodeSum) & Yes & No & Maybe & Supervised baseline; missing data \\
Yun et al.\ 2024 (P-CodeSum) & Yes & Partial & Likely & Most promising runnable baseline \\
\bottomrule
\end{tabular}
\end{table*}
  % Uncomment when integrating

\begin{table}[ht]
\centering
\caption{Dataset Composition by Repository}
\label{tab:dataset-composition}
\begin{tabular}{lcc}
\toprule
Repository & Functions & Proportion \\
\midrule
requests & 5 & 33.3\% \\
pandas & 5 & 33.3\% \\
scikit-learn & 5 & 33.3\% \\
\midrule
\textbf{Total} & \textbf{15} & \textbf{100.0\%} \\
\bottomrule
\end{tabular}
\end{table}

\subsection{Prompting Methods}
We compare two prompting strategies:

\begin{enumerate}
  \item \textbf{Generic Prompting:} A zero-shot prompt containing only a
    summarization instruction, the target function code, and a warning not to
    hallucinate unsupported behavior.
  \item \textbf{Project-Specific ICL:} A few-shot prompt containing 2--3
    example function--summary pairs from the same repository as the target
    function, followed by the target function code and the same instructions.
\end{enumerate}

\subsection{Evaluation}
Each of the 30~generated summaries (15~per method) was manually labeled as:
\begin{itemize}
  \item \textbf{Correct:} The summary accurately describes the function's
    behavior.
  \item \textbf{Partial:} The summary is partially accurate but omits
    important behavior or includes minor inaccuracies.
  \item \textbf{Incorrect:} The summary contains a material error or
    describes behavior the function does not exhibit.
\end{itemize}

% ============================================================
\section{Results}
\label{sec:results}

Table~\ref{tab:result-summary} presents the evaluation results.

\begin{table}[ht]
\centering
\caption{Summary of Evaluation Results}
\label{tab:result-summary}
\begin{tabular}{lccccc}
\toprule
Method & Correct & Partial & Incorrect & Total & Correctness Rate \\
\midrule
Generic Prompting & 7 & 4 & 4 & 15 & 46.7\% \\
Project-Specific ICL & 11 & 3 & 1 & 15 & 73.3\% \\
\bottomrule
\end{tabular}
\end{table}

Project-specific ICL achieved a correctness rate of 73.3\%, compared to 46.7\%
for generic prompting---an improvement of 26.6~percentage points. The
incorrect-summary rate dropped from 26.7\% (4/15) to 6.7\% (1/15).

\subsection{What Was Seen}
Generic prompting tended to produce summaries that were either too vague
(leading to Partial labels) or that confused similar-sounding operations
(leading to Incorrect labels). For example, the generic summary for
\texttt{normalize} confused normalization with standardization, and the summary
for \texttt{unique} incorrectly described a set-based implementation.

Project-specific ICL summaries were more aligned with each project's vocabulary
and conventions. The ICL examples appeared to help the LLM distinguish between
closely related operations (e.g., normalization vs.\ standardization in
scikit-learn).

% ============================================================
\section{Discussion}
\label{sec:discussion}

\subsection{Results and Implications}
The improvement from 46.7\% to 73.3\% correctness suggests that project-specific
context provides meaningful grounding for LLM-generated summaries. The reduction
in incorrect summaries (from 4 to 1) is particularly encouraging, as incorrect
summaries pose a greater risk than partial ones in practice.

These results are consistent with prior findings by Ahmed and
Devanbu~\cite{ahmed2022fewshot} and Yun et al.~\cite{yun2024pcodesum}, who
also found that project-specific context improves code summarization quality.

\subsection{Threats to Validity}
\label{sec:validity}

\textbf{Internal Validity.}
Manual labeling is inherently subjective. Different annotators may assign
different labels to the same summary. We do not report inter-rater reliability
because the study used a single annotator. Future work should employ multiple
annotators and report agreement metrics.

\textbf{External Validity.}
The study covers only 15~functions from 3~Python repositories. Results may not
generalize to other languages, projects, or function types. The small sample
size precludes robust statistical significance testing.

\textbf{Construct Validity.}
The three-level labeling scheme (Correct, Partial, Incorrect) is coarse.
Finer-grained evaluation metrics (e.g., semantic similarity scores) could
provide additional insight.

\textbf{Reproducibility.}
Generated summaries depend on the specific LLM, its version, and API
parameters at the time of generation. Results may differ with different
models or API versions. All artifacts are provided to support replication
of the evaluation, if not the generation step.

% ============================================================
\section{Baseline Reproducibility}
\label{sec:baseline-reproducibility}

We attempted to reproduce three prior systems referenced in our study:

\begin{enumerate}
  \item \textbf{Ahmed and Devanbu 2022~\cite{ahmed2022fewshot}:}
    Repository at
    \url{https://github.com/toufiqueparag/few_shot_code_summarization}.
    Repo cloned and environment set up successfully, but the script failed
    at the API call because the original model \texttt{code-davinci-002}
    is deprecated. \textbf{Status: failed unmodified.}
  \item \textbf{Ahmad et al.\ 2020 (NeuralCodeSum)~\cite{ahmad2020neuralcodesum}:}
    Repository at \url{https://github.com/wasiahmad/NeuralCodeSum}.
    Repo cloned, dependencies repaired, pipeline launched, but failed because
    required Java dataset files were missing and the bundled download path
    did not produce a valid archive. \textbf{Status: failed unmodified.}
  \item \textbf{Yun et al.\ 2024 (P-CodeSum)~\cite{yun2024pcodesum}:}
    Repository at
    \url{https://github.com/Linshuhuai/P-CodeSum}.
    Repo installed successfully and scripts launched, but default
    \texttt{train.sh}/\texttt{evaluate.sh} point to nonexistent data paths;
    the repo contains alternate data directories and log files were created
    successfully. \textbf{Status: partial; most promising runnable baseline.}
\end{enumerate}

Our direct empirical baseline is Generic Prompting. Prior systems are treated
as reproducibility reference baselines and are documented in the artifact
repository.

% ============================================================
\section{Future Work}
\label{sec:future-work}

\begin{itemize}
  \item Scale the study to more functions, repositories, and programming
    languages.
  \item Employ multiple annotators and report inter-rater agreement.
  \item Compare against automated metrics (BLEU, ROUGE, BERTScore).
  \item Investigate how the number and quality of ICL examples affect
    summary reliability.
  \item Evaluate with multiple LLMs to assess generalizability across models.
\end{itemize}

% ============================================================
\section{Conclusion}
\label{sec:conclusion}

In this exploratory study, project-specific in-context learning improved the
correctness rate of LLM-generated code summaries from 46.7\% to 73.3\% and
reduced the incorrect-summary rate from 26.7\% to 6.7\%. While the small
sample size limits statistical generalizability, the results suggest that
grounding LLM prompts in project-specific context is a promising direction
for improving the reliability of automated code summarization. All artifacts
are publicly available for replication.

% ============================================================
\section*{Replication Artifacts}

All data, prompts, generated outputs, evaluation labels, analysis scripts,
literature review artifacts (100-paper citation list, knee analysis, thematic
classification, Venn counts), and baseline reproduction notes are available at:

\begin{center}
% TODO: Replace <username> with your actual GitHub username
\url{https://github.com/<username>/project-specific-code-summary-icl}
\end{center}

To reproduce the result tables:
\begin{verbatim}
pip install -r requirements.txt
python scripts/validate_artifacts.py
python scripts/compute_results.py
python scripts/summarize_labels.py
python scripts/generate_latex_tables.py
\end{verbatim}

% ============================================================
\bibliographystyle{IEEEtran}
\bibliography{references}

\end{document}
